\documentclass[11pt,a4paper]{article}

% =========================================================
% Packages
% =========================================================
\usepackage[utf8]{inputenc}
\usepackage[T1]{fontenc}
\usepackage{lmodern}
\usepackage{geometry}
\usepackage{amsmath, amssymb, amsthm, mathtools}
\usepackage{hyperref}
\usepackage{enumitem}
\usepackage{setspace}
\usepackage{titlesec}
\usepackage{booktabs}
\usepackage{array}
\usepackage{longtable}

\geometry{margin=1in}
\onehalfspacing

\titleformat{\section}{\large\bfseries}{\thesection.}{0.5em}{}
\titleformat{\subsection}{\normalsize\bfseries}{\thesubsection.}{0.5em}{}
\titleformat{\subsubsection}{\normalsize\itshape}{\thesubsubsection.}{0.5em}{}

% =========================================================
% Theorem-like environments
% =========================================================
\newtheorem{definition}{Definition}
\newtheorem{principle}{Principle}
\newtheorem{proposition}{Proposition}
\newtheorem{remark}{Remark}

% =========================================================
% Commands
% =========================================================
\newcommand{\At}{A_t}
\newcommand{\Atn}{A_{t+1}}
\newcommand{\DeltaA}{\Delta_{\Omega}A_t}
\newcommand{\OmegaT}{\Omega_t}
\newcommand{\OmegaD}{\Omega_{\Delta}}
\newcommand{\OmegaE}{\Omega_E}
\newcommand{\ST}{S_t}
\newcommand{\UT}{U_t}
\newcommand{\MT}{M_t}
\newcommand{\Wct}{Wc_t}
\newcommand{\Dt}{D_t}
\newcommand{\Toxt}{Tox_t}
\newcommand{\Tt}{T_t}
\newcommand{\Ht}{H_t}
\newcommand{\It}{I_t}
\newcommand{\Cap}{\operatorname{Cap}}
\newcommand{\Valid}{\operatorname{Valid}}
\newcommand{\Resource}{\operatorname{Resource}}
\newcommand{\dist}{\operatorname{d}}
\newcommand{\argmax}{\operatorname*{arg,max}}

% =========================================================
% Title
% =========================================================
\title{\textbf{A Formal Equation of Existence:}\
Survival, Development, Differentiating Ground, and Resource Transformation}

\author{Kevin T.N}

\date{}

\begin{document}

\maketitle

\begin{abstract}
This paper proposes a formal framework for understanding existence as a dynamic system process rather than a static condition. The central claim is that any system that exists must be determinable within a differentiating ground, denoted by $\Omega$, in which the system $A$ can be located, bounded, and distinguished from at least one non-$A$. Within this ground, the system survives by selectively accessing, filtering, transforming, and expelling non-$A$ components. Development is defined as the differential by which the system moves from its current survival state $A_t$ to a new survival state $A_{t+1}$. The primitive equation is
[
A_{t+1}=A_t+\Delta_{\Omega}A_t,
]
where
[
\Delta_{\Omega}A_t=A_{t+1}-A_t.
]
The paper develops this into a general survival function that includes usable resource extraction, maintenance, repair, waste cost, dissipation, toxic accumulation, temporal mismatch, homeostatic regulation, and structural information continuity. The expanded equation is
[
A_{t+1}
=======

A_t
+
f\left[
\int_{\Omega_t \setminus A_t}
V_A(x,t)\alpha_A(x,t)P_A(x,t)\eta_A(x,t)k_A(x,t)G_A(x,t),dx
-----------------------------------------------------------

## M_t(A_t)

## Wc_t(A_t)

## D_t(A_t)

## Tox_t(A_t)

T_t(A_t)
\right],
]
subject to boundary, homeostatic, core-continuity, and information-continuity constraints. The framework aims to provide a general formal language for discussing existence across physical, chemical, biological, cognitive, organizational, and social systems.
\end{abstract}

\textbf{Keywords:} existence, survival, development, systems theory, resource transformation, differentiating ground, boundary, entropy, homeostasis, information, waste, identity

\section{Introduction}

Existence is often treated as a static predicate. Something exists if it is present, detectable, or not absent. While this minimal meaning may be useful in ordinary language, it does not explain how a system continues to exist through time. A biological organism, a cell, a mind, a company, a society, or a technological system does not merely exist by being present. It persists by continuously maintaining itself through the intake, filtering, conversion, repair, regulation, and expulsion of what lies outside it.

This paper proposes a formal equation of existence. The goal is not to reduce all forms of existence to a single mechanical formula, but to build a general framework that identifies the minimum structural conditions under which a system may be said to survive, develop, decay, or collapse. The framework begins from a simple observation: a system cannot be determined in absolute isolation. For a system $A$ to be identifiable, there must be a domain in which $A$ can be distinguished from what is not $A$. This domain is denoted by $\Omega$ and is called a differentiating ground.

Once a system is located in a differentiating ground, its existence is not a passive fact. It must maintain a boundary, access non-$A$ components, convert some of them into usable resources, expel or neutralize what it cannot convert, compensate for dissipation, regulate internal variables, repair damage, and preserve enough structural information to remain itself across transformation. Development, in this framework, is not external to survival. It is the differential by which one survival state becomes a wider, stronger, or more capable survival state.

The primitive equation is:
[
A' = A + (A' - A).
]
In temporal form:
[
A_{t+1}=A_t+\Delta_{\Omega}A_t,
]
where:
[
\Delta_{\Omega}A_t=A_{t+1}-A_t.
]

Here, $A_t$ is the current survival state of the system. $A_{t+1}$ is the next survival state after transformation. The differential $\Delta_{\Omega}A_t$ is development, but only if it increases the system's capacity without breaking its core continuity.

The paper proceeds as follows. Section 2 defines the differentiating ground $\Omega$. Section 3 distinguishes between the minimal differential ground and the full existential ground. Section 4 defines the system $A$ and its boundary. Section 5 defines non-$A$ and resource. Section 6 formulates usable resource extraction. Section 7 introduces cost, waste, dissipation, toxicity, and temporal mismatch. Section 8 introduces homeostasis, feedback, and information continuity. Section 9 defines survival, development, stagnation, decline, and collapse. Section 10 presents the master equation. Section 11 distinguishes optimization from selection. Section 12 interprets the framework across physical, chemical, biological, cognitive, organizational, and social domains. The paper concludes by summarizing existence as the recursive process by which a system converts what is not itself into the means of continuing to be itself.

\section{The Differentiating Ground}

\subsection{The Need for a Ground}

A system cannot be determined in an absolute void. In order for $A$ to be identified, there must be a domain in which $A$ can be located and distinguished from at least one non-$A$. This domain is not necessarily physical space. It may be a physical environment, a chemical reaction space, an ecological niche, a cognitive field, a linguistic system, a market, a social structure, a mathematical state space, or a space of possibilities.

\begin{definition}[Differentiating Ground]
A differentiating ground $\Omega$ is any determinate domain in which a system $A$ can be located and distinguished from at least one non-$A$.
\end{definition}

The minimal validity condition is:
[
\Valid(\Omega,A)
\Longleftrightarrow
A\in\Omega
\wedge
\exists X\in\Omega:X\neq A.
]

This means that a ground is valid for $A$ if and only if $A$ is in the ground and there exists at least one distinguishable component in the same ground that is not $A$.

Without such a ground, $A$ has no determinate boundary. Without a boundary, there is no inside and outside. Without inside and outside, there is no input, output, waste, access, conversion, or development. Therefore, existence does not begin merely with $A$; it begins with the condition under which $A$ can be distinguished from non-$A$.

\subsection{The Boundary Condition}

A boundary is not merely a spatial edge. It is a distinction between what belongs to the system and what does not. Therefore, the boundary of $A$ is possible only within a differentiating ground.

Let $B(A)$ denote the boundary of $A$. Then:
[
B(A) \text{ is meaningful only if } \exists X\in\Omega:X\neq A.
]

The boundary separates the system from non-systemic components, but it also mediates exchange. In simple systems, boundary may function primarily as distinction. In living and adaptive systems, boundary is selective: it regulates what enters, what is blocked, what exits, and what is retained.

\section{Minimal Differential Ground and Full Existential Ground}

\subsection{The Minimal Differential Ground}

The smallest possible ground for difference contains $A$ and one distinguishable other. In the most reduced case, this other may be the system's own possible next state $A'$:
[
\Omega_{\Delta}(A)={A,A'},\quad A'\neq A.
]

This ground is sufficient to define difference:
[
\Delta_{\Omega}A=A'-A.
]

In temporal form:
[
\Delta_{\Omega}A_t=A_{t+1}-A_t.
]

Thus, $A'$ may function as a minimal non-$A$. It allows $A$ to be compared with another state. This is enough to define change, but it is not enough to define systemic existence.

\subsection{The Limitation of the Binary Ground}

If the entire ground is only:
[
\Omega={A,A'},
]
then the system is locked into a closed binary relation. It can distinguish $A$ from $A'$, but it cannot determine what $A$ consumes, what $A$ expels, what constrains $A$, what alternatives exist, what feedback is available, or whether $A'$ is viable.

The binary ground allows the statement:
\begin{quote}
$A$ is different from $A'$.
\end{quote}

It does not answer:
\begin{quote}
What does $A$ live on? What does it convert? What does it reject? What limits it? What stabilizes $A'$? Is $A'$ a development, a deformation, or a collapse?
\end{quote}

Therefore, $A'$ is sufficient for minimal difference but insufficient for full existence.

\subsection{The Full Existential Ground}

A full existential ground must contain a richer field of non-$A$ components:
[
\Omega_E(A)={A,A',X_1,X_2,\ldots,X_n},
]
where at least some $X_i$ are neither $A$ nor $A'$.

These components may function as resources, constraints, waste receivers, competitors, collaborators, signals, feedback channels, tools, alternative states, or destabilizing forces. Therefore:
[
\Omega_{\Delta}(A)\subset\Omega_E(A).
]

\begin{principle}[Differential Ground and Existential Ground]
The minimal differential ground is sufficient for defining change. The full existential ground is required for defining survival, development, waste, constraint, feedback, and stabilization.
\end{principle}

In short, $A'$ is enough for $A$ to know that it can be different. It is not enough for $A$ to know what it survives on.

\section{The System}

\subsection{System Composition}

Let a system at time $t$ be denoted by:
[
A_t.
]

The system may be represented as:
[
A_t={E_t,L_t,K_t,B_t,I_t,q_t},
]
where:
\begin{itemize}[leftmargin=2em]
\item $E_t$ denotes the elements of the system;
\item $L_t$ denotes the relations among elements;
\item $K_t$ denotes the core structure of the system;
\item $B_t$ denotes the boundary of the system;
\item $I_t$ denotes structural information, memory, or code;
\item $q_t$ denotes the vector of internal state variables.
\end{itemize}

A system is not merely a set of elements. It is a structured relation among elements. The relations are not free. They require maintenance. A biological organism must maintain tissues, organs, and regulatory pathways. An organization must maintain communication, coordination, procedures, trust, and decision structures. A society must maintain institutions, laws, infrastructure, language, and symbolic order.

Thus, survival is not only the maintenance of elements. It is also the maintenance of relations.

\subsection{Open Systems}

The systems most relevant to this framework are open systems. They exchange energy, matter, information, or signals with their surroundings. They do not preserve themselves by remaining sealed. They preserve themselves by selective openness.

A fully closed system may retain a boundary, but it cannot sustain development if it cannot receive new resources or expel waste. Conversely, a fully open system with no selective boundary loses its internal organization. Therefore, survival requires neither total closure nor total openness, but regulated exchange.

\section{Non-A and Resource}

\subsection{Non-A}

Within a differentiating ground $\Omega$, any component that is not $A$ is non-$A$:
[
X\in\Omega\setminus A.
]

A non-$A$ component may be a resource, a toxin, a constraint, a signal, a competitor, a collaborator, a waste receiver, a tool, a possible future state, or a source of disturbance. Therefore, non-$A$ is not automatically useful. It is simply what $A$ can be distinguished from.

\subsection{Resource as Convertible Non-A}

A resource is not an absolute category. A component becomes a resource for $A$ only if it can be accessed, admitted through the boundary, and transformed into something useful for survival or development.

\begin{definition}[Resource]
A component $x$ is a resource for $A$ if $x$ is a non-$A$ component in $\Omega$ that $A$ can access, admit through its boundary, and convert into usable value.
\end{definition}

Formally:
[
\Resource_A(x)
\Longleftrightarrow
x\in\Omega_t\setminus A_t
\wedge
\alpha_A(x,t)>0
\wedge
P_A(x,t)>0
\wedge
G_A(x,t)>0
\wedge
\eta_A(x,t)>0.
]

Here:
\begin{itemize}[leftmargin=2em]
\item $\alpha_A(x,t)$ is the accessibility of $x$ to $A$;
\item $P_A(x,t)$ is the permeability or boundary permission of $x$;
\item $G_A(x,t)$ is the free-energy potential or transformation gradient of $x$ relative to $A$;
\item $\eta_A(x,t)$ is the conversion efficiency of $x$ by $A$.
\end{itemize}

A thing may exist in the same ground as $A$ but not be a resource. It may be inaccessible, blocked by the boundary, energetically unusable, kinetically too slow to transform, or destructive if absorbed.

\section{Access, Boundary Permeability, and Conversion}

\subsection{Access}

Access is the ability of $A$ to reach or interact with $x$:
[
\alpha_A(x,t).
]

Access may depend on distance, position, technology, permission, recognition, mobility, institutional right, cognitive attention, or communication channel. In physical systems, access may be spatial. In economic systems, it may be legal or logistical. In cognitive systems, it may be attentional or conceptual.

\subsection{Boundary Permeability}

Permeability is the degree to which $x$ can pass through the selective boundary of $A$:
[
P_A(x,t).
]

Access and permeability are distinct. A system may encounter a component but reject it. A cell may encounter a molecule but not allow it through the membrane. A person may encounter information but filter it out. An organization may encounter talent, capital, or data but fail to admit it into its structure.

Thus:
[
\alpha_A(x,t)>0 \nRightarrow P_A(x,t)>0.
]

\subsection{Conversion Capacity}

Conversion is the ability of $A$ to transform $x$ into usable structure, energy, function, information, or capacity. In simplified form, this may be represented by:
[
C_A(x,t).
]

However, conversion may be decomposed into several components:
[
C_A(x,t)\sim \eta_A(x,t)k_A(x,t)G_A(x,t).
]

Here:
\begin{itemize}[leftmargin=2em]
\item $\eta_A(x,t)$ is efficiency;
\item $k_A(x,t)$ is kinetic rate;
\item $G_A(x,t)$ is free-energy potential or gradient.
\end{itemize}

This decomposition reflects a basic constraint: something may be valuable in principle but useless in practice if it cannot be converted quickly enough, efficiently enough, or under available energetic conditions.

\subsection{Kinetics and Catalysis}

The conversion rate $k_A(x,t)$ may depend on barriers and catalytic support. A general abstract expression is:
[
k_A(x,t)=k_0\cdot \chi_A(x,t)\cdot e^{-\lambda_A(x,t)},
]
where:
\begin{itemize}[leftmargin=2em]
\item $k_0$ is a baseline rate;
\item $\chi_A(x,t)$ represents catalytic support;
\item $\lambda_A(x,t)$ represents the transformation barrier.
\end{itemize}

The exact functional form may vary by domain. The key point is that conversion is not merely a yes-or-no condition. It has rate, barrier, and facilitation.

\section{The Usable Resource Function}

The usable resource available to $A$ within $\Omega_t$ is the portion of non-$A$ that can be accessed, admitted, and converted.

For a continuous ground:
[
U_t(A_t|\Omega_t)
=================

\int_{\Omega_t\setminus A_t}
V_A(x,t)\alpha_A(x,t)P_A(x,t)\eta_A(x,t)k_A(x,t)G_A(x,t),dx.
]

For a discrete ground:
[
U_t(A_t|\Omega_t)
=================

\sum_{X_i\in\Omega_t,,X_i\neq A_t}
V_A(X_i,t)\alpha_A(X_i,t)P_A(X_i,t)\eta_A(X_i,t)k_A(X_i,t)G_A(X_i,t).
]

Here:
\begin{itemize}[leftmargin=2em]
\item $V_A(x,t)$ is the potential value of $x$ for $A$;
\item $\alpha_A(x,t)$ is access;
\item $P_A(x,t)$ is boundary permeability;
\item $\eta_A(x,t)$ is efficiency;
\item $k_A(x,t)$ is conversion rate;
\item $G_A(x,t)$ is transformation gradient.
\end{itemize}

This equation is central. It states that the environment is not automatically resource. Usable resource is produced by the relation between system, boundary, access, conversion, and ground.

\section{Maintenance, Repair, Waste, and Dissipation}

\subsection{Maintenance Cost}

Let $M_t(A_t)$ denote the maintenance cost of $A_t$. This cost includes the cost of maintaining elements, relations, and repair mechanisms:
[
M_t(A_t)=M_E(t)+M_L(t)+M_R(t),
]
where:
\begin{itemize}[leftmargin=2em]
\item $M_E(t)$ is the cost of maintaining elements;
\item $M_L(t)$ is the cost of maintaining relations;
\item $M_R(t)$ is the cost of repair.
\end{itemize}

A system may collapse not because it lacks parts, but because it can no longer maintain the relations among its parts. In organisms, this may appear as regulatory failure. In organizations, it may appear as coordination cost. In societies, it may appear as institutional breakdown.

\subsection{Waste}

Let $W_t(A)$ denote the waste produced by $A$. Waste is the portion of transformed or encountered non-$A$ that the system cannot use, cannot retain, or must expel.

Waste is relational. What is waste for $A$ may be resource for another system $B$:
[
W_A=R_B
]
if $B$ has the conversion capacity that $A$ lacks.

\begin{principle}[Relative Waste]
Waste is not an absolute category. Waste is resource outside the conversion capacity or retention capacity of a given system.
\end{principle}

\subsection{Waste Cost and Toxic Accumulation}

Let $Wc_t(A)$ denote the cost of processing, neutralizing, transporting, or expelling waste. Let $\beta_A(W,t)$ denote the waste export capacity of $A$.

If waste production exceeds export capacity, waste accumulates as toxicity:
[
Tox_t(A)=\max(0,W_t(A)-\beta_A(W,t)).
]

Thus, waste is not merely a cost. It may become an internal destabilizing force. A system may collapse not because it lacks input, but because it cannot expel output.

\subsection{Dissipation}

Let $D_t(A)$ denote dissipation cost. Dissipation includes losses due to heat, friction, delay, disorder, noise, entropy production, or any domain-equivalent form of irreversible loss.

No conversion is perfectly free. Therefore, survival must compensate not only for maintenance and waste, but also for dissipation:
[
D_t(A)>0
]
for any non-trivial transformation in real systems.

\subsection{Temporal Mismatch}

Let $T_t(A)$ denote temporal mismatch cost. This represents the cost produced when the rhythm of input, conversion, output, repair, or stabilization is misaligned.

A system can fail even when total resources are sufficient if:
\begin{itemize}[leftmargin=2em]
\item resources arrive too slowly;
\item resources arrive too quickly;
\item conversion cannot keep up with input;
\item waste production exceeds export rhythm;
\item damage occurs faster than repair;
\item adaptation occurs slower than environmental change.
\end{itemize}

Time is therefore not merely an index. It is a rate constraint on transformation.

\section{Homeostasis, Feedback, and Information}

\subsection{Homeostatic State}

Let:
[
q_t(A)
]
denote the vector of internal state variables of $A$. These may include physical, chemical, biological, cognitive, organizational, or social variables depending on the system.

Let:
[
q_A^*
]
denote the viable range or target region for these variables.

The homeostatic deviation is:
[
H_t(A)=d(q_t(A),q_A^*).
]

A system remains internally viable only if:
[
H_t(A)\leq h,
]
where $h$ is the maximum tolerable deviation.

This condition is crucial. A system may have sufficient resources and still collapse if its internal variables leave the viable region.

\subsection{Feedback and Regulation}

Let:
[
F_t(A)
]
denote the feedback or regulatory function of $A$. Feedback regulates access, boundary permeability, conversion, repair, waste export, and internal variables.

Feedback may:
\begin{itemize}[leftmargin=2em]
\item increase intake when resources are low;
\item decrease intake when overload occurs;
\item increase waste export when toxicity accumulates;
\item activate repair when damage occurs;
\item adjust conversion pathways;
\item stabilize internal variables;
\item alter boundary permeability;
\item update strategy based on environmental signals.
\end{itemize}

Without feedback, a system is merely a passive flow structure. With feedback, it becomes adaptive.

\subsection{Structural Information}

Let:
[
I_t(A)
]
denote the structural information, memory, or code that allows $A$ to reproduce its organization across transformation. This may include genetic information, regulatory architecture, memory, habits, procedures, symbolic order, or institutional rules, depending on the domain.

Structural information evolves as:
[
I_{t+1}(A)=I_t(A)+\Delta I_t-\varepsilon_I,
]
where:
\begin{itemize}[leftmargin=2em]
\item $\Delta I_t$ is newly acquired or updated structural information;
\item $\varepsilon_I$ is information loss, distortion, noise, or error.
\end{itemize}

Information continuity requires:
[
d_I[I_{t+1}(A),I_t(A)]\leq \theta_I.
]

Energy allows a system to operate. Information allows a system to reproduce its organization as itself.

\section{The Survival Function}

The survival level of $A_t$ within $\Omega_t$ is:
[
S_t(A_t|\Omega_t)
=================

## U_t(A_t|\Omega_t)

## M_t(A_t)

## Wc_t(A_t)

## D_t(A_t)

## Tox_t(A_t)

T_t(A_t).
]

Substituting the usable resource function:
[
S_t(A_t|\Omega_t)
=================

\int_{\Omega_t\setminus A_t}
V_A(x,t)\alpha_A(x,t)P_A(x,t)\eta_A(x,t)k_A(x,t)G_A(x,t),dx
-----------------------------------------------------------

## M_t(A_t)

## Wc_t(A_t)

## D_t(A_t)

## Tox_t(A_t)

T_t(A_t).
]

\begin{definition}[Survival Function]
The survival function of a system $A_t$ in a differentiating ground $\Omega_t$ is the net viability produced by usable resource extraction minus maintenance, repair, waste, dissipation, toxicity, and temporal mismatch costs.
\end{definition}

The minimal energetic-material survival condition is:
[
S_t(A_t|\Omega_t)\geq 0.
]

But full survival requires more than a non-negative resource balance. It also requires homeostatic stability, core continuity, and information continuity:
[
Survival
\Longleftrightarrow
S_t(A_t|\Omega_t)\geq 0
\wedge
H_t(A_t)\leq h
\wedge
d_{\Omega}[K(A_{t+1}),K(A_t)]\leq\theta
\wedge
d_I[I_{t+1}(A),I_t(A)]\leq\theta_I.
]

This definition prevents a common error: treating survival as mere energy surplus. A system can have surplus resources and still collapse if its internal variables, core structure, or structural information break.

\section{Development}

Development is the differential between the current survival state and the next survival state:
[
\Delta_{\Omega}A_t=A_{t+1}-A_t.
]

Thus:
[
A_{t+1}=A_t+\Delta_{\Omega}A_t.
]

However, not every difference is development. Difference may be neutral, destructive, regressive, or transformative. Development requires survival plus increased capacity.

Let:
[
\Cap(A_t)
]
denote the capacity of $A_t$. Capacity includes the ability to access non-$A$, filter through boundary, convert resources, regulate internal variables, export waste, repair damage, preserve information, and stabilize new states.

Development occurs when:
[
Development
\Longleftrightarrow
Survival
\wedge
\Cap(A_{t+1})>\Cap(A_t)
\wedge
d_{\Omega}[K(A_{t+1}),K(A_t)]\leq\theta.
]

A more expansive condition may include:
[
\Omega_{t+1}\succeq\Omega_t,
]
where $\Omega_{t+1}\succeq\Omega_t$ means that the new ground is not necessarily physically larger, but richer in determination, access, conversion, feedback, alternatives, or stabilizing pathways.

\begin{principle}[Development as Expanded Survival]
Development is the process by which a system absorbs a differential change into a new survival state with greater capacity while preserving core continuity.
\end{principle}

Development is therefore not mere growth. Growth may overload the system. Expansion may destabilize it. Development requires that the new state be stabilizable as a new survival state.

\section{Decay, Stagnation, and Collapse}

The framework allows several system states to be distinguished.

\subsection{Stagnation}

A system stagnates when it survives but does not increase capacity:
[
S_t(A|\Omega)\geq 0
\wedge
\Cap(A_{t+1})=\Cap(A_t).
]

Stagnation is not immediate collapse. It is survival without capacity expansion.

\subsection{Decline}

A system declines when it remains viable in the short term but loses capacity:
[
S_t(A|\Omega)\geq 0
\wedge
\Cap(A_{t+1})<\Cap(A_t).
]

Decline may precede collapse if the loss of capacity eventually makes survival impossible.

\subsection{Collapse}

A system collapses when at least one critical condition fails:
[
Collapse
\Longleftrightarrow
S_t(A|\Omega)<0
\vee
H_t(A)>h
\vee
d_{\Omega}[K(A_{t+1}),K(A_t)]>\theta
\vee
d_I[I_{t+1},I_t]>\theta_I.
]

Collapse may therefore result from resource deficit, internal instability, core rupture, or information loss.

\section{The Master Equation}

The core equation of existence is:
[
\boxed{
A_{t+1}=A_t+f(S_t(A_t|\Omega_t))
}
]

The expanded survival function is:
[
\boxed{
S_t(A_t|\Omega_t)
=================

\int_{\Omega_t\setminus A_t}
V_A(x,t)\alpha_A(x,t)P_A(x,t)\eta_A(x,t)k_A(x,t)G_A(x,t),dx
-----------------------------------------------------------

## M_t(A_t)

## Wc_t(A_t)

## D_t(A_t)

## Tox_t(A_t)

T_t(A_t)
}
]

Therefore, the expanded equation of existence is:
[
\boxed{
A_{t+1}
=======

A_t
+
f\left[
\int_{\Omega_t\setminus A_t}
V_A(x,t)\alpha_A(x,t)P_A(x,t)\eta_A(x,t)k_A(x,t)G_A(x,t),dx
-----------------------------------------------------------

## M_t(A_t)

## Wc_t(A_t)

## D_t(A_t)

## Tox_t(A_t)

T_t(A_t)
\right]
}
]

Subject to:
[
\boxed{
A_t\in\Omega_t
\wedge
\exists X\in\Omega_t:X\neq A_t
}
]

[
\boxed{
H_t(A_t)\leq h
}
]

[
\boxed{
d_{\Omega}[K(A_{t+1}),K(A_t)]\leq\theta
}
]

[
\boxed{
d_I[I_{t+1}(A),I_t(A)]\leq\theta_I
}
]

This equation expresses existence as a constrained transformation. The system does not simply become $A_{t+1}$. It becomes $A_{t+1}$ through resource extraction, boundary filtering, conversion, dissipation, waste handling, regulation, repair, and information continuity.

\section{Optimization and Selection}

\subsection{Regulated Optimization}

For systems capable of directed regulation, planning, or control, the next state may be modeled as an optimization problem:
[
A_{t+1}
=======

\argmax_A S_t(A|\Omega),
]
subject to:
[
H_t(A)\leq h,
]
[
d_{\Omega}[K(A_{t+1}),K(A_t)]\leq\theta,
]
[
d_I[I_{t+1},I_t]\leq\theta_I.
]

This form is most appropriate for agents, organizations, institutions, designed systems, and cognitive systems that can represent alternatives and select among them.

\subsection{Selection Without Central Intention}

In natural evolutionary contexts, it is often misleading to say that a system consciously maximizes survival. Biological evolution operates through variation, inheritance, and differential survival-reproduction success across populations.

Let $A_i$ denote a variant in a population. Then:
[
P(A_i,t+1)\propto P(A_i,t)\cdot F(A_i|\Omega_t),
]
where $F(A_i|\Omega_t)$ is the fitness of variant $A_i$ within $\Omega_t$.

This distinction is important. Regulated systems may optimize. Natural populations are filtered by selection.

\begin{principle}[Optimization and Selection]
Directed systems may optimize survival under constraints. Undirected natural populations are selected through differential fitness across variants.
\end{principle}

\section{Replication and Lineage Existence}

The equation $A_t\rightarrow A_{t+1}$ describes individual continuity. Some systems also maintain themselves through replication, reproduction, or transmission.

Let:
[
Rep_t(A)
]
denote the replication capacity of $A$.

Individual existence may be represented as:
[
A_t\rightarrow A_{t+1}.
]

Lineage existence may be represented as:
[
A_t\rightarrow{A_{t+1}^{(1)},A_{t+1}^{(2)},\ldots,A_{t+1}^{(n)}}.
]

Lineage continuity requires that offspring or successors preserve sufficient core and structural information:
[
d_{\Omega}[K(A_{t+1}^{(i)}),K(A_t)]\leq\theta_r,
]
and:
[
d_I[I(A_{t+1}^{(i)}),I(A_t)]\leq\theta_{Ir}.
]

This applies not only to biological reproduction, but also to organizational scaling, cultural transmission, software replication, education, institutional reproduction, and memetic continuity.

\section{Interpretive Domains}

\subsection{Physical Systems}

In physical systems, $\Omega$ is the environment or surrounding domain in which the system is distinguished. $A$ is the system, $B(A)$ its boundary, $U_t$ usable energy or work potential, and $D_t$ dissipation. The framework aligns with the idea that a system transforms through exchange and cannot generate usable capacity without energetic constraint. The existence equation therefore includes dissipation to prevent the false assumption of lossless transformation.

\subsection{Chemical Systems}

In chemical systems, non-$A$ components may be reactants, catalysts, solvents, ions, or environmental conditions. The variables $G_A$, $k_A$, and $\eta_A$ become especially important. A substance may be present but chemically irrelevant if there is no favorable gradient, if the reaction is too slow, or if the system lacks the catalytic pathway to transform it.

The chemical interpretation clarifies that resource is not mere presence. A component becomes resource only when transformation is possible under the given conditions.

\subsection{Biological Systems}

In biological systems, the boundary becomes a selective membrane or equivalent regulatory surface. $P_A$ corresponds to selective admission. $U_t$ includes usable energy and matter. $H_t$ corresponds to homeostatic stability. $I_t$ corresponds to genetic, epigenetic, regulatory, and structural information. $F_t$ corresponds to feedback regulation. $Rep_t$ corresponds to reproductive or replicative capacity.

Biological existence is therefore not reducible to energy intake. It requires selective boundaries, metabolism, waste export, repair, regulation, information preservation, and reproduction.

\subsection{Cognitive Systems}

In cognitive systems, $\Omega$ may be an informational, experiential, linguistic, or social field. Non-$A$ components include signals, concepts, memories, contradictions, other minds, and possible future selves. Access may be attention. Boundary permeability may be filtering or relevance selection. Conversion is understanding. Waste may be noise, confusion, contradiction, trauma, or unprocessed information. Structural information is memory and self-model continuity.

A mind survives cognitively when it can process incoming reality without losing coherence. It develops when it can process more complex variables while preserving identity and internal stability.

\subsection{Organizations}

In organizations, $\Omega$ may be a market, institutional environment, technological platform, or regulatory field. Non-$A$ components include capital, labor, customers, data, suppliers, competitors, tools, and regulations. Boundary permeability includes hiring, procurement, access control, culture, and onboarding. Conversion is operational capacity. Waste includes bureaucracy, failed projects, debt, internal conflict, and reputational damage.

An organization survives when usable resources exceed maintenance, coordination, waste, and repair costs. It develops when it increases capacity without breaking its core operating logic.

\subsection{Societies}

In societies, $\Omega$ includes territory, ecology, population, institutions, language, law, technology, economy, culture, and external systems. Non-$A$ components include resources, threats, opportunities, contradictions, external powers, innovations, and waste fields. Structural information includes law, language, historical memory, symbolic order, institutional pattern, and education.

A society survives when it can maintain coherence while processing resources, waste, conflict, and change. It develops when it expands its capacity to process complexity without destroying the continuity that makes it recognizable as itself.

\section{Discussion}

The proposed framework makes several contributions.

First, it shifts existence from a static predicate to a dynamic process. To exist is not merely to be present. To exist is to maintain continuity through transformation.

Second, it establishes the differentiating ground as a prior condition. There is no determinate $A$ without $\Omega$ and non-$A$. This avoids treating systems as abstract isolated objects without boundary, relation, or context.

Third, it distinguishes the minimal ground of difference from the full ground of existence. The pair ${A,A'}$ is enough for difference, but not enough for survival. Full existence requires resources, constraints, waste, feedback, and stabilizing relations.

Fourth, it defines resource relationally. Resource is not what exists in the environment. Resource is what the system can access, admit, and convert.

Fifth, it integrates boundary as a selective mechanism. A system survives not by being open to everything, but by regulating what enters and exits.

Sixth, it treats waste and toxicity as central, not secondary. A system may collapse due to output failure even when input is abundant.

Seventh, it incorporates dissipation and temporal mismatch. Transformation is costly and rate-limited.

Eighth, it introduces homeostasis and information continuity as survival constraints. Resource surplus alone is insufficient. The system must keep its internal variables within viable bounds and preserve enough structural information to remain itself.

Finally, it clarifies development. Development is not growth in size, nor mere change. Development is the successful absorption of a differential into a new survival state with greater capacity and preserved core continuity.

\section{Conclusion}

This paper has proposed a formal equation of existence based on differentiating ground, survival, development, boundary, resource transformation, waste, dissipation, homeostasis, feedback, and information continuity.

The primitive form is:
[
A_{t+1}=A_t+\Delta_{\Omega}A_t.
]

The differential is:
[
\Delta_{\Omega}A_t=A_{t+1}-A_t.
]

The survival function is:
[
S_t(A_t|\Omega_t)
=================

## U_t(A_t|\Omega_t)

## M_t(A_t)

## Wc_t(A_t)

## D_t(A_t)

## Tox_t(A_t)

T_t(A_t).
]

The usable resource function is:
[
U_t(A_t|\Omega_t)
=================

\int_{\Omega_t\setminus A_t}
V_A(x,t)\alpha_A(x,t)P_A(x,t)\eta_A(x,t)k_A(x,t)G_A(x,t),dx.
]

Therefore:
[
A_{t+1}=A_t+f(S_t(A_t|\Omega_t)).
]

The full equation is:
[
A_{t+1}
=======

A_t
+
f\left[
\int_{\Omega_t\setminus A_t}
V_A(x,t)\alpha_A(x,t)P_A(x,t)\eta_A(x,t)k_A(x,t)G_A(x,t),dx
-----------------------------------------------------------

## M_t(A_t)

## Wc_t(A_t)

## D_t(A_t)

## Tox_t(A_t)

T_t(A_t)
\right].
]

Subject to:
[
A_t\in\Omega_t,
\quad
\exists X\in\Omega_t:X\neq A_t,
]
[
H_t(A_t)\leq h,
]
[
d_{\Omega}[K(A_{t+1}),K(A_t)]\leq\theta,
]
[
d_I[I_{t+1}(A),I_t(A)]\leq\theta_I.
]

The core thesis may be stated as follows:

\begin{quote}
Existence is the recursive process by which a system, located within a differentiating ground, selectively converts what is not itself into the means of continuing to be itself, while producing a differential that may become a new survival state.
\end{quote}

A system survives when it can maintain its core through regulated exchange. It develops when it can transform a differential into greater capacity without losing itself. It collapses when usable resources fail, internal variables leave viable bounds, waste accumulates beyond export capacity, dissipation cannot be compensated, structural information breaks, or the core exceeds its threshold of continuity.

In the most compact form:

\begin{quote}
A system exists when it can digest non-$A$ in order to remain $A$. A system develops when it can digest more of non-$A$ without ceasing to be $A$.
\end{quote}

\appendix

\section{Symbol Table}

\begin{longtable}{p{0.18\textwidth}p{0.32\textwidth}p{0.42\textwidth}}
\toprule
\textbf{Symbol} & \textbf{Name} & \textbf{Meaning} \
\midrule
\endhead

$\Omega_t$ & Differentiating ground & Domain in which $A_t$ can be located and distinguished from non-$A$. \

$\Omega_{\Delta}$ & Minimal differential ground & Minimal ground containing $A$ and $A'$; sufficient for difference. \

$\Omega_E$ & Full existential ground & Richer ground containing resources, constraints, waste relations, feedback, and alternatives. \

$A_t$ & System state & System at time $t$. \

$A_{t+1}$ & Next system state & New survival state after transformation. \

$\Delta_{\Omega}A_t$ & Developmental differential & Difference between $A_{t+1}$ and $A_t$ within $\Omega$. \

$E_t$ & Elements & Parts of the system. \

$L_t$ & Relations & Relations among elements. \

$K_t$ & Core & Core structure or continuity of the system. \

$B_t$ & Boundary & Distinction and selective interface between $A$ and non-$A$. \

$I_t$ & Structural information & Memory, code, or organizational information preserving the system. \

$q_t$ & Internal state vector & Internal variables that must remain within viable range. \

$X$ & Non-$A$ component & Any component in $\Omega$ that is not $A$. \

$U_t$ & Usable resource & Extractable and convertible value available to $A$. \

$V_A$ & Potential value & Potential value of $x$ for $A$. \

$\alpha_A$ & Access & Ability of $A$ to reach or interact with $x$. \

$P_A$ & Permeability & Ability of $x$ to pass through the boundary of $A$. \

$\eta_A$ & Efficiency & Conversion efficiency. \

$k_A$ & Kinetic rate & Rate of conversion. \

$G_A$ & Gradient & Free-energy potential or transformation gradient. \

$M_t$ & Maintenance cost & Cost of maintaining elements, relations, and repair. \

$W_t$ & Waste & Output that the system cannot use or must expel. \

$Wc_t$ & Waste cost & Cost of handling or neutralizing waste. \

$\beta_A$ & Waste export capacity & Capacity to expel waste. \

$Tox_t$ & Toxic accumulation & Waste accumulation beyond export capacity. \

$D_t$ & Dissipation cost & Irreversible loss due to transformation. \

$T_t$ & Temporal mismatch cost & Cost of rate mismatch between input, conversion, output, repair, and stabilization. \

$H_t$ & Homeostatic deviation & Distance between internal state and viable range. \

$F_t$ & Feedback & Regulatory function. \

$\Cap(A_t)$ & Capacity & System's ability to access, filter, convert, regulate, repair, expel, preserve information, and stabilize. \

$Rep_t(A)$ & Replication capacity & Ability to reproduce or transmit the system's core/information into successors. \

\bottomrule
\end{longtable}

\section{Compact Formal Summary}

[
\Valid(\Omega,A)
\Longleftrightarrow
A\in\Omega
\wedge
\exists X\in\Omega:X\neq A
]

[
\Omega_{\Delta}(A)={A,A'},\quad A'\neq A
]

[
\Omega_{\Delta}(A)\subset\Omega_E(A)
]

[
A_{t+1}=A_t+\Delta_{\Omega}A_t
]

[
\Delta_{\Omega}A_t=A_{t+1}-A_t
]

[
\Resource_A(x)
\Longleftrightarrow
x\in\Omega_t\setminus A_t
\wedge
\alpha_A(x,t)>0
\wedge
P_A(x,t)>0
\wedge
G_A(x,t)>0
\wedge
\eta_A(x,t)>0
]

[
U_t(A_t|\Omega_t)
=================

\int_{\Omega_t\setminus A_t}
V_A(x,t)\alpha_A(x,t)P_A(x,t)\eta_A(x,t)k_A(x,t)G_A(x,t),dx
]

[
M_t(A_t)=M_E(t)+M_L(t)+M_R(t)
]

[
Tox_t(A)=\max(0,W_t(A)-\beta_A(W,t))
]

[
H_t(A)=d(q_t(A),q_A^*)
]

[
I_{t+1}(A)=I_t(A)+\Delta I_t-\varepsilon_I
]

[
S_t(A_t|\Omega_t)
=================

## U_t(A_t|\Omega_t)

## M_t(A_t)

## Wc_t(A_t)

## D_t(A_t)

## Tox_t(A_t)

T_t(A_t)
]

[
Survival
\Longleftrightarrow
S_t(A_t|\Omega_t)\geq0
\wedge
H_t(A_t)\leq h
\wedge
d_{\Omega}[K(A_{t+1}),K(A_t)]\leq\theta
\wedge
d_I[I_{t+1}(A),I_t(A)]\leq\theta_I
]

[
Development
\Longleftrightarrow
Survival
\wedge
\Cap(A_{t+1})>\Cap(A_t)
\wedge
d_{\Omega}[K(A_{t+1}),K(A_t)]\leq\theta
]

[
Collapse
\Longleftrightarrow
S_t(A_t|\Omega_t)<0
\vee
H_t(A_t)>h
\vee
d_{\Omega}[K(A_{t+1}),K(A_t)]>\theta
\vee
d_I[I_{t+1},I_t]>\theta_I
]

[
A_{t+1}=A_t+f(S_t(A_t|\Omega_t))
]

[
A_{t+1}
=======

A_t
+
f\left[
\int_{\Omega_t\setminus A_t}
V_A(x,t)\alpha_A(x,t)P_A(x,t)\eta_A(x,t)k_A(x,t)G_A(x,t),dx
-----------------------------------------------------------

## M_t(A_t)

## Wc_t(A_t)

## D_t(A_t)

## Tox_t(A_t)

T_t(A_t)
\right]
]

\section{References}

\begin{thebibliography}{99}

\bibitem{bertalanffy1968}
von Bertalanffy, L. (1968).
\textit{General System Theory: Foundations, Development, Applications}.
George Braziller.

\bibitem{ashby1956}
Ashby, W. R. (1956).
\textit{An Introduction to Cybernetics}.
Chapman & Hall.

\bibitem{wiener1948}
Wiener, N. (1948).
\textit{Cybernetics: Or Control and Communication in the Animal and the Machine}.
MIT Press.

\bibitem{shannon1948}
Shannon, C. E. (1948).
A mathematical theory of communication.
\textit{Bell System Technical Journal}, 27, 379--423, 623--656.

\bibitem{schrodinger1944}
Schrödinger, E. (1944).
\textit{What Is Life? The Physical Aspect of the Living Cell}.
Cambridge University Press.

\bibitem{prigogine1984}
Prigogine, I., & Stengers, I. (1984).
\textit{Order Out of Chaos: Man's New Dialogue with Nature}.
Bantam Books.

\bibitem{maturana1980}
Maturana, H. R., & Varela, F. J. (1980).
\textit{Autopoiesis and Cognition: The Realization of the Living}.
D. Reidel Publishing.

\bibitem{kauffman1993}
Kauffman, S. A. (1993).
\textit{The Origins of Order: Self-Organization and Selection in Evolution}.
Oxford University Press.

\bibitem{maynardsmith1999}
Maynard Smith, J., & Szathmáry, E. (1999).
\textit{The Origins of Life: From the Birth of Life to the Origin of Language}.
Oxford University Press.

\bibitem{odum1983}
Odum, H. T. (1983).
\textit{Systems Ecology: An Introduction}.
Wiley.

\bibitem{nicolis1977}
Nicolis, G., & Prigogine, I. (1977).
\textit{Self-Organization in Nonequilibrium Systems}.
Wiley.

\bibitem{friston2010}
Friston, K. (2010).
The free-energy principle: a unified brain theory?
\textit{Nature Reviews Neuroscience}, 11, 127--138.

\end{thebibliography}

\end{document}
